% Appendix B: Symbolic computation of the equations of motion (~1600 w body prose
% + TOML figure + B.5 equations + deferred-derived-fields footnote).
% Per genre_conventions_software.md §(i). Renamed from "Symbolic stage"
% — the BHL convention is content-bearing names (compare PSALTer v2's
% "Construction of operators", "Moore--Penrose pseudoinversion") rather
% than pipeline-stage labels.

\section{Symbolic computation of the equations of motion}%
\label{SymbolicComputation}

From a Lagrangian, a background ansatz, and a parameter set, the
symbolic computation produces a JSON specification of the linearised
field equations and the Hamiltonian density.
\cref{Architecture} sets out the high-level pipeline, with the data-flow
diagram shown in \cref{ArchitectureDiagram};
the main-text overview is \cref{ComputationalApproach}, and the
equation outputs for the surveyed theories are collected in
\cref{SymbolicOutputs}.
Every partial differential equation evolved by the numerical solver of
\cref{NumericalEvolution} is the output of this pipeline; for
higher-curvature theories, an additional perturbative-reduction stage
described in \cref{ConstraintBarrier} precedes the variation.

\paragraph*{xAct environment}
The symbolic stage runs in \Mathematica{}~\cite{mathematica} on the
\xAct{}~\cite{martingarcia2007xact} suite, which provides the
abstract-index tensor algebra, automatic component decomposition, and
perturbation theory.
Three \xAct{} packages are loaded: \texttt{xTensor} for abstract-index
algebra and the functional derivative \texttt{VarD};
\xPert{}~\cite{brizuela2009xpert} for metric perturbation theory through
\texttt{DefTensorPerturbation}; and \xTras{}~\cite{nutma2014xtras} for
the Einstein--Cartan operators that include the torsion-inclusive
covariant derivative used throughout the PGT theories of
\cref{PGTBackground}.
Two further \xAct{} packages, \texttt{xCoba} for charts and
\texttt{xPerm} for index symmetries, enter as transitive dependencies of
the first three.
The manifold, metric, Levi-Civita and Riemann--Cartan covariant
derivatives, and the torsion tensor are declared once at session start;
the Christoffel components, the inverse metric, and the Levi-Civita
tensor are computed once and cached against the declared metric.

\paragraph*{Lagrangian input from TOML}
The theory is specified by a single TOML file --- a human-readable
configuration format used widely in scientific software --- which
isolates everything theory-dependent from the derivation pipeline so
the same code can serve any quadratic Lagrangian without per-theory
editing.
The TOML declares the spacetime (dimension, signature, coordinate
names); the dynamical fields, with their tensor type and index
symmetries; any algebraically-derived tensors such as the
electromagnetic field strength $\FieldF{_\mu_\nu}$; the background
fields; the parameter symbol set; the Lagrangian expression in
\xAct{} syntax; and optional blocks for gauge fixing, linearisation,
and perturbative reduction.
A representative theory file --- the Einstein--Maxwell setup that
underlies the Gertsenshtein analysis of \cref{Results} --- is shown in
\cref{TOMLExample}.
The pipeline as a whole is the map
\begin{align*}
  \textsc{TOML}\!\bigl[\mathcal{L},\,\{\varphi_i\},\,\{p_a\}\bigr]
  \;\longmapsto\;
  \textsc{JSON}\!\bigl[\{\textsf{EOM}_i\},\,\mathcal{H}\bigr],
\end{align*}
whose internal stages occupy the rest of this appendix.

\begin{figure}[t]
\begin{lstlisting}[basicstyle=\ttfamily\footnotesize,frame=single,
    framerule=0.4pt,columns=fullflexible,
    keepspaces=true,xleftmargin=1em]
[spacetime]
dimension = 4
signature = "minkowski"
coordinates = ["t", "x", "y", "z"]

[[fields]]
name = "h"           # metric perturbation
type = "tensor"
rank = 2
symmetry = "symmetric"

[[fields]]
name = "A"           # electromagnetic potential
type = "vector"

[[derived_fields]]
name = "F"           # F_{ab} = d A
type = "tensor"
rank = 2
symmetry = "antisymmetric"
definition = "CD[-a][A[-b]] - CD[-b][A[-a]]"

[[background_fields]]
name  = "Abar"       # uniform magnetic field along x
type  = "vector"
components = ["0", "0", "-B0 * z[]", "0"]

[parameters]
names  = ["B0", "kappa"]

[lagrangian]
expression = "-(1/4) F[-a,-b] F[a,b] + (1/kappa^2) RicciScalarCD[]"

[[linearization.matter_perturbations]]
field             = "A"
perturbation_name = "a"
background        = "Abar"

[gauge]
field   = "A"
preset  = "lorenz"
\end{lstlisting}
\caption{Representative \TIDAL{} input file: the Einstein--Maxwell
theory on a uniform background magnetic field. The file declares the
spacetime, the dynamical fields (metric perturbation, electromagnetic
potential), the algebraically-derived field strength, the background
magnetic vector, the parameter set, the Lagrangian in \xAct{} syntax,
the matter-perturbation block that drives the multi-field linearisation
of \cref{Linearisation}, and the Lorenz gauge. All other theories
in the surveys of \cref{Results} share this structure and differ only
in their declared fields and Lagrangian.}
\label{TOMLExample}
\end{figure}

\paragraph*{Linearisation around backgrounds}
The pipeline implements the background-and-perturbation decomposition
of \cref{Linearisation} by expanding the action to second order in the
perturbations using \xPert{}~\cite{brizuela2009xpert}, with each
dynamical field perturbed independently around its declared background.
The multi-field generalisation perturbs the metric and any declared
matter fields simultaneously.

For each field the perturbation is registered as a tensor in \xAct{}:
the metric carries one through \texttt{SetupMetricPerturbation}, and
each matter field $\varphi^{\mathcal{I}}$ declared in the
matter-perturbations block of \cref{TOMLExample} carries one through
\texttt{DefTensorPerturbation}, yielding
\begin{align}
  g_{\mu\nu} = \eta_{\mu\nu} + \varepsilon\,h_{\mu\nu},
  \qquad
  \varphi^{\mathcal{I}} = \bar{\varphi}^{\mathcal{I}}
                          + \varepsilon\,\delta\varphi^{\mathcal{I}},
  \label{MultiFieldPerturbation}
\end{align}
where $\mathcal{I}$ stands for the field's tensor indices, suppressed
throughout this paragraph; \cref{MultiFieldPerturbation} is the
pipeline-side reading of \cref{PerturbationAnsatz}.
Backgrounds are supplied through the background-fields block of
\cref{TOMLExample}, with their component values fixed by
\texttt{xCoba}'s \texttt{ComponentValue} declarations --- for example
the uniform magnetic background of \cref{PerturbationAnsatz},
$\bar{A}_y = -\Bzero\,z$.\footnote{Antisymmetric derived tensors whose
definition references a perturbation field --- most prominently
$\FieldF{_\mu_\nu} = \rcD{_\mu}\FieldA{_\nu} - \rcD{_\nu}\FieldA{_\mu}$
when $\FieldA{}$ is itself a matter perturbation --- must remain
abstract through the perturbation expansion and be substituted only
afterwards. Without this deferral, the canonicalisation step that
follows collapses the antisymmetric kinetic structure to zero before
perturbation indices are assigned, and the photon equations emerge as
non-dynamical constraints. The pipeline detects deferral candidates
automatically from the TOML symmetry annotation and threads them
through the expansion as abstract tensors with their own
\texttt{DefTensorPerturbation} registration.}

We expand the action to second order in the perturbation parameter
$\varepsilon$ by applying \texttt{Perturbation} of order two to
$\sqrt{-g}\,\mathcal{L}$ followed by \texttt{ExpandPerturbation}.
The resulting $\mathcal{L}^{(2)}$ contains two kinds of terms: products
of first-order perturbations (the \textsf{LI}[1] symbols of \xPert{}),
which constitute the kinetic and coupling structure of the linearised
theory, and linear couplings to isolated second-order perturbation
tensors (the \textsf{LI}[2] symbols, parametrising the
$\mathcal{O}(\varepsilon^{2})$ corrections to the fields themselves).
The first kind we keep verbatim; the second we drop via the
substitution rule $\textsf{LI}[2] \mapsto 0$, since the linearised
equations of motion are obtained by variation with respect to
first-order fields and are unaffected by the \textsf{LI}[2] content.
The retained
\begin{align}
  \mathcal{L}^{(2)}
  \;=\;
  \textsf{ExpandPerturbation}\!\bigl[
    \textsf{Perturbation}[\sqrt{-g}\,\mathcal{L},\,2]
  \bigr]\Big|_{\textsf{LI}[2]\,\mapsto\,0}
  \label{QuadraticExpansion}
\end{align}
is the linearised action of the multi-field system.

The remaining abstract covariant derivatives of background tensors are
resolved against the \texttt{ComponentValue} declarations during the
chart conversion, so $\rcD{_z}\bar{A}_y = -\Bzero$ for the uniform
background and a coordinate-dependent expression for a localised one
--- the mechanism by which position-dependent coupling coefficients
enter the JSON specification without special-casing in the solver.

\paragraph*{Euler--Lagrange derivation}
The field equation for every dynamical perturbation is obtained by
variational differentiation of the quadratic action $\mathcal{L}^{(2)}$
built in \cref{QuadraticExpansion}: for fields $\{\varphi_i\}$ of arbitrary
tensor type, the equation of motion is
\begin{align}
  \frac{\partial \mathcal{L}^{(2)}}{\partial \varphi_i}
  \;-\; \rcD{_\mu}\!\frac{\partial \mathcal{L}^{(2)}}{\partial(\rcD{_\mu}\varphi_i)}
  \;+\; \rcD{_\mu}\rcD{_\nu}\!\frac{\partial \mathcal{L}^{(2)}}{\partial(\rcD{_\mu}\rcD{_\nu}\varphi_i)}
  \;-\; \cdots \;=\; 0,
  \label{EulerLagrange}
\end{align}
where the series terminates at the highest derivative order present in
$\mathcal{L}^{(2)}$.
This is the only stage of the pipeline that consumes the Lagrangian as
a single object; every later stage operates on the equations themselves.
The differentiation is performed by \xAct{}'s
\texttt{VarD}~\cite{martingarcia2007xact}, which handles the chain rule
and integration-by-parts for arbitrary tensor structure and arbitrary
derivative chains, and returns the equation in abstract-index form
before any component reduction.

The equations are kept in second-order velocity form, with explicit
velocity slots $v_i = \partial_t\varphi_i$ alongside the fields; the
rationale --- the kinetic-matrix definition, the canonical-momentum
inversion failures that the velocity form sidesteps, and the
consequences for solver dispatch --- is the subject of
\cref{NumericalEvolution}'s \emph{Velocity form versus canonical phase
space} block.
For Lagrangians containing higher-curvature terms multiplied by a
small parameter, which generate equations of motion of order higher
than two in time, the pipeline applies a perturbative substitution
before the variation; the algorithm, its scope, and the structural
obstruction that limits it are deferred to \cref{ConstraintBarrier}.

\paragraph*{Component decomposition}
The abstract-index equations returned by \texttt{VarD} are reduced to
scalar component equations on the declared chart before numerical
evaluation.
Four operations are performed in sequence: the metric is separated
from any contracted indices; the expression is converted to chart
components by \texttt{ToBasis}, which projects an abstract-index
tensor expression onto the chart's coordinate basis,
\begin{align}
  \textsf{ToBasis}\!\bigl[T_{\mu_1\cdots\mu_n}\bigr]
  \;=\; T_{i_1\cdots i_n}\,
        \mathrm{d}x^{i_1} \otimes \cdots \otimes \mathrm{d}x^{i_n},
  \label{ToBasisProjection}
\end{align}
with Roman indices labelling chart components; repeated indices are
summed by \texttt{TraceBasisDummy}; and the abstract covariant
derivative is converted to partial derivatives plus the appropriate
connection corrections,
\begin{align}
  \rcD{_\mu}\varphi_{\nu}
  \;=\;
  \PD{_\mu}\varphi_{\nu}
  \;-\; \bigl(\rCon{^\sigma_\mu_\nu}
              + K^{\sigma}{}_{\mu\nu}\bigr)\,\varphi_{\sigma},
  \label{CovariantDerivativeExpansion}
\end{align}
where $\mathring{\Gamma}$ is the Levi-Civita part and $K$ the
contortion; for tensors of higher rank the routine emits one such
connection term per slot, drawn from the geometric cache established
at session start.
The cache specialises to vanishing Christoffels whenever the declared
metric matrix has no coordinate dependence, and to vanishing contortion
whenever the theory has no torsion, so the routine produces the
simplest expression compatible with the declared geometry.
Canonicalisation of the full quadratic Lagrangian is performed in a
single pass rather than by batching across additive terms, because
batched canonicalisation can collapse antisymmetric kinetic structure
that spans multiple summands; an \texttt{Expand}-only fallback is used
when the single pass exceeds the canonicalisation timeout, with
per-component canonicalisation in the chart-conversion step recovering
the canonical form downstream.

\paragraph*{Gauge fixing}%
\label{GaugeFixing}
When a gauge is to be fixed --- an option, not a requirement, since
the surveys of \cref{Results} run at both fixed and unfixed gauges
depending on what is known about a theory's symmetry structure ---
\TIDAL{} offers two structurally distinct mechanisms: some gauges are
most naturally imposed before the Euler--Lagrange variation, others
after.

The first mechanism adds a gauge-fixing term to the Lagrangian before
variation~\cite{faddeev1967feynman}.
For a vector field $\FieldA{_\mu}$, the Lorenz gauge is realised by
the augmentation
\begin{align}
  \mathcal{L} \;\longrightarrow\;
  \mathcal{L} \;-\; \frac{1}{2\xi}\bigl(\rcD{_\mu}\FieldA{^\mu}\bigr)^{2},
  \label{LorenzGauge}
\end{align}
with the gauge parameter $\xi$ controlling the weight of the residual
symmetry; the augmented Lagrangian then enters the variational
pipeline in the usual way.
The second mechanism imposes a constraint on the field components
after decomposition, as in the transverse-traceless prescription that
sets $h_{0\mu} = 0$ and $h^\mu{}_\mu = 0$ on the metric-perturbation
components --- a constraint expressible only at the component level
because it refers to component indices, not abstract ones.

Six named presets cover the gauges in routine use: Lorenz for vector
fields, de Donder for the metric perturbation, Coulomb for the spatial
part of a vector field, temporal ($A_0 = 0$), axial ($A_3 = 0$), and
transverse-traceless.
Arbitrary user-specified gauges are declared by writing the gauge
expression directly in the \texttt{[[gauge]]} TOML block, which is
treated as either a Lagrangian term or a component constraint
according to a flag.
The choice of gauge interacts with the constraint elimination performed
at the numerical-evolution stage: gauge-fixed kinetic operators
sometimes still produce equations of zeroth time-derivative order,
which the JSON-export step identifies and which the constraint
pre-solve of \cref{NumericalEvolution} handles before the system is
integrated.

\paragraph*{JSON export}
The symbolic stage terminates by writing a JSON specification, the
sole interface to the numerical side of \TIDAL{}; fixing this interface
cleanly is what allows parameter sweeps and grid changes without
re-derivation.
The specification has five top-level entries: a spacetime block
(dimension, signature, coordinate names); the field list, with each
field's time-derivative order recorded explicitly --- two for
wave-like, one for first-order, zero for constraint; the equations,
each carried as a linear combination of named differential operators
acting on the dynamical fields with parameter-dependent coefficients;
the canonical block, which holds the Hamiltonian density in the same
operator format and is used for energy diagnostics; and a metadata
block recording the source Lagrangian, the gauge, the linearisation
flag, the parameter values, and a derivation hash.
The operator alphabet --- identity, full and directional Laplacians,
directional gradients, mixed cross-derivatives, a first-order time
derivative, the biharmonic, and a generic mixed-derivative form --- is
fixed and small.
Each term carries both its numerical coefficient and its symbolic
expression, so a runtime parameter override never regenerates the
JSON; the equation specification is the centre arrow of
\cref{ArchitectureDiagram}.

\paragraph*{Failure modes and diagnostics}
The symbolic stage has three classes of failure that the pipeline
detects and surfaces as derivation artefacts rather than as silent
numerical pathologies downstream.
The canonicalisation of the full quadratic Lagrangian is super-linear
in term count, and the default six-hundred-second ceiling triggers
the \texttt{Expand}-only fallback discussed above --- equivalent on
linear theories, and cheaper.
Malformed Lagrangians, such as free-index miscounts or contractions
across symbolic dummies that \xAct{} cannot resolve, are caught by
\xAct{}'s built-in index-validation pass and reported with the
offending term.
A post-export validator confirms that every operator and field
reference is known and that the canonical block matches the kinetic
structure implied by the field list's time-derivative orders; failures
point the user to the responsible Lagrangian term.
Surfacing failures here, rather than in the numerical stage, is what
made the surveys of \cref{Results} tractable; the complementary
numerical-validation pass of \cref{Validation} catches the remainder.
